\documentclass[11pt]{article}
\usepackage{amsmath, amssymb, amsthm}
\usepackage{mathtools}
\usepackage{bm}

\title{Theoretical Interpretation of a Lifted Product Code Implementation}
\author{Debashis Saikia}
\date{}

\begin{document}

\maketitle

\section{Overview}

The lifted product construction defines a chain complex over a group algebra

\begin{equation}
R = \mathbb{F}_2[G]
\end{equation}

where $G$ is a finite group. Given matrices

\begin{equation}
A \in R^{m_A \times n_A}, \quad B \in R^{m_B \times n_B},
\end{equation}

we construct a 3-term chain complex:

\begin{equation}
R^{n_A n_B}
\overset{\partial_2}{\longrightarrow}
R^{m_A n_B} \oplus R^{n_A m_B}
\overset{\partial_1}{\longrightarrow}
R^{m_A m_B}.
\end{equation}

The boundary maps are:

\begin{equation}
\partial_2 =
\begin{bmatrix}
A \otimes I \
I \otimes B
\end{bmatrix},
\quad
\partial_1 =
\begin{bmatrix}
I \otimes B & A \otimes I
\end{bmatrix}.
\end{equation}

The CSS code is defined as:

\begin{equation}
H_X = \partial_1, \quad H_Z = \partial_2^{T}.
\end{equation}

\section{Finite Group Representation}

\subsection{Code}
\begin{verbatim}
class FiniteGroup:
def **init**(self, mul_table):
self.mul = mul_table
self.n = len(mul_table)
\end{verbatim}

\subsection{Theory}

This defines a finite group

\begin{equation}
G = {g_0, \dots, g_{L-1}}
\end{equation}

with multiplication:

\begin{equation}
g_i \cdot g_j = g_{\text{mul}[i][j]}.
\end{equation}

The integer $L = |G|$ is stored as:

\begin{equation}
L = \texttt{self.n}.
\end{equation}

\section{Group Algebra Elements}

\subsection{Code}
\begin{verbatim}
class GroupAlgebraElement:
def **init**(self, coeffs, group):
self.coeffs = np.array(coeffs) % 2
\end{verbatim}

\subsection{Theory}

Each element represents:

\begin{equation}
a = \sum_{g \in G} a_g g, \quad a_g \in \mathbb{F}_2.
\end{equation}

The vector representation is:

\begin{equation}
(a_{g_0}, a_{g_1}, \dots, a_{g_{L-1}}).
\end{equation}

\subsection{Addition}

\begin{verbatim}
return self.coeffs ^ other.coeffs
\end{verbatim}

\begin{equation}
(a + b)_g = a_g + b_g \mod 2.
\end{equation}

\subsection{Multiplication}

\begin{verbatim}
k = G.mul[i][j]
result[k] ^= 1
\end{verbatim}

This implements:

\begin{equation}
\left(\sum a_g g\right)\left(\sum b_h h\right)
=\sum_{g,h} a_g b_h (g h),
\end{equation}

which matches the group algebra multiplication:

\begin{equation}
ab = \sum_{x \in G} \left( \sum_{gh = x} a_g b_h \right) x.
\end{equation}

\section{Left and Right Module Actions}

\begin{verbatim}
def right_mul(self, other):
def left_mul(self, other):
\end{verbatim}

\subsection{Theory}

We distinguish:
\begin{itemize}
\item Right multiplication: $a \cdot b$
\item Left multiplication: $b \cdot a$
\end{itemize}

This is essential because:
\begin{itemize}
\item $A$ is a \textbf{right} $R$-module
\item $B$ is a \textbf{left} $R$-module
\end{itemize}

\section{Matrices over the Group Algebra}

\begin{verbatim}
class GroupAlgebraMatrix:
def **init**(self, data):
\end{verbatim}

\subsection{Theory}

Represents matrices:

\begin{equation}
A \in R^{m_A \times n_A}, \quad B \in R^{m_B \times n_B}.
\end{equation}

Each entry lies in $\mathbb{F}_2[G]$.

\section{Regular Representations}

\subsection{Right Regular Representation}

\begin{verbatim}
def right_regular_matrix(ga):
\end{verbatim}

\begin{equation}
\rho_R(g): e_i \mapsto e_{i g}.
\end{equation}

\begin{equation}
M_{j,i} = 1 \quad \text{if } j = i g.
\end{equation}

\subsection{Left Regular Representation}

\begin{verbatim}
def left_regular_matrix(ga):
\end{verbatim}

\begin{equation}
\rho_L(g): e_i \mapsto e_{g i}.
\end{equation}

\section{Lifting to Binary Matrices}

\begin{verbatim}
def lift_matrix(GA, side="right"):
\end{verbatim}

\subsection{Theory}

Each element of $\mathbb{F}_2[G]$ is replaced by its matrix representation:

\begin{equation}
\mathbb{F}_2[G] \cong \mathbb{F}_2^{L \times L}.
\end{equation}

Thus:

\begin{equation}
A \mapsto \widetilde{A} \in \mathbb{F}_2^{(mL) \times (nL)}.
\end{equation}

\section{Tensor Product Construction}

\subsection{Code}
\begin{verbatim}
A_kron = np.kron(HA, IB)
B_kron = np.kron(IA, HB)
\end{verbatim}

\subsection{Theory}

These correspond to:

\begin{equation}
A \otimes I, \quad I \otimes B.
\end{equation}

\section{Boundary Operator $\partial_2$}

\begin{verbatim}
self.d2 = np.vstack([A_kron, B_kron])
\end{verbatim}

\begin{equation}
\partial_2 =
\begin{bmatrix}
A \otimes I \
I \otimes B
\end{bmatrix}.
\end{equation}

\section{Boundary Operator $\partial_1$}

\begin{verbatim}
left_block  = np.kron(ImA, HB)
right_block = np.kron(HA, ImB)
\end{verbatim}

\begin{equation}
\partial_1 =
\begin{bmatrix}
I \otimes B & A \otimes I
\end{bmatrix}.
\end{equation}

\section{CSS Code Construction}

\begin{verbatim}
self.HX = self.d1
self.HZ = self.d2.T
\end{verbatim}

\begin{equation}
H_X = \partial_1, \quad H_Z = \partial_2^T.
\end{equation}

\section{Chain Complex Condition}

\begin{verbatim}
HX @ HZ.T == 0
\end{verbatim}

\begin{equation}
\partial_1 \circ \partial_2 = 0.
\end{equation}

\section{Code Parameters}

\subsection{Length}

\begin{equation}
n = n_A n_B |G|.
\end{equation}

\subsection{Dimension}

\begin{verbatim}
k = n - rank(HX) - rank(HZ)
\end{verbatim}

\begin{equation}
k = \dim \ker \partial_1 - \dim \operatorname{im} \partial_2.
\end{equation}

\section{Conclusion}

Each component of the implementation corresponds exactly to the lifted product construction:
\begin{itemize}
\item Group algebra $\mathbb{F}_2[G]$
\item Tensor product over $R$
\item Boundary maps $\partial_1, \partial_2$
\item CSS construction
\end{itemize}

Thus the code is a faithful computational realization of the lifted product framework.

\end{document}
